\documentclass[a4paper,11pt]{article}
\usepackage{amsmath,amssymb,epsf,epsfig,times}
\usepackage[all]{xy}
\usepackage{color}
\usepackage{subfigure}
\usepackage{url,cite}
\newtheorem{theorem}{Theorem}[section]
\newtheorem{lemma}{Lemma}[section]
\def\proof{\noindent{\it Proof: }}
\def\QED{\mbox{\rule[0pt]{1.5ex}{1.5ex}}}
\def\endproof{\hspace*{\fill}~\QED\par\endtrivlist\unskip}
\newcommand{\re}{\mathbb{R}}
\def\sV{\mathcal{V}}
\def\sS{\mathcal{S}}
\def\sQ{\mathcal{Q}}

\newcommand{\mc}{\mbox{: }}

\newcommand{\normsq}[1]{\left\|#1\right\|^2}
\newcommand{\norm}[1]{\left\|#1\right\|}
%\newcommand{\sgn}[1]{\mbox{sgn}(#1)}
\newcommand{\pde}[2]{\frac{\partial #1}{\partial #2}}
\newcommand{\fundef}[3]{#1:#2\to #3}
\newcommand{\abs}[1]{\left|#1\right|}
\newcommand{\mymatrix}[2]{\left(\begin{array}{#1}#2\end{array}\right)}
\newcommand{\defeq}{\stackrel{\triangle}{=}}
\newcommand{\paren}[1]{\left(#1\right)}
%\theoremstyle{plain} \newtheorem{theorem}{Theorem}
%\theoremstyle{plain} \newtheorem{algorithm}{Algorithm}
\newtheorem{axiom}[theorem]{Axiom}
\newtheorem{definition}[theorem]{Definition}
\newtheorem{assumption}[theorem]{Assumption}
\newtheorem{example}[theorem]{Example}
%\theoremstyle{plain}\newtheorem{lemma}{Lemma}
%\newtheorem{proposition}[theorem]{Proposition}
\newtheorem{remark}[theorem]{Remark}
\newtheorem{corollary}[theorem]{Corollary}

\newcommand{\Acal}{\mathcal{A}}
\newcommand{\Bcal}{\mathcal{B}}
\newcommand{\Ccal}{\mathcal{C}}
\newcommand{\Dcal}{\mathcal{D}}
\newcommand{\Ecal}{\mathcal{E}}
\newcommand{\Fcal}{\mathcal{F}}
\newcommand{\Gcal}{\mathcal{G}}
\newcommand{\Hcal}{\mathcal{H}}
\newcommand{\Ical}{\mathcal{I}}
\newcommand{\Jcal}{\mathcal{J}}
\newcommand{\Kcal}{\mathcal{K}}
\newcommand{\Lcal}{\mathcal{L}}
\newcommand{\Mcal}{\mathcal{M}}
\newcommand{\Ncal}{\mathcal{N}}
\newcommand{\Ocal}{\mathcal{O}}
\newcommand{\Pcal}{\mathcal{P}}
\newcommand{\Qcal}{\mathcal{Q}}
\newcommand{\Rcal}{\mathcal{R}}
\newcommand{\Scal}{\mathcal{S}}
\newcommand{\Tcal}{\mathcal{T}}
\newcommand{\Ucal}{\mathcal{U}}
\newcommand{\Vcal}{\mathcal{V}}
\newcommand{\Wcal}{\mathcal{W}}
\newcommand{\Xcal}{\mathcal{X}}
\newcommand{\Ycal}{\mathcal{Y}}
\newcommand{\Zcal}{\mathcal{Z}}


\def\omegavec{\boldsymbol{\omega}}
\newcommand{\alphabf}{\boldsymbol{\alpha}}
\newcommand{\omegabf}{\boldsymbol{\omega}}
\def\omegavec{\boldsymbol{\omega}}
\newcommand{\taubf}{\boldsymbol{\tau}}
\newcommand{\qbf}{\mathbf{q}}
\newcommand{\ybf}{\mathbf{y}}
\newcommand{\pbf}{\mathbf{p}}
\newcommand{\rbf}{\mathbf{r}}
\newcommand{\ebf}{\mathbf{e}}
\newcommand{\onebf}{\mathbf{1}}
\newcommand{\zerobf}{\mathbf{0}}
\newcommand{\abf}{\mathbf{a}}
\newcommand{\ibf}{\mathbf{i}}
\newcommand{\jbf}{\mathbf{j}}
\newcommand{\kbf}{\mathbf{k}}
\newcommand{\vbf}{\mathbf{v}}
\newcommand{\wbf}{\mathbf{\omega}}
\newcommand{\fbf}{\mathbf{f}}
\newcommand{\zbf}{\mathbf{z}}
\newcommand{\xbf}{\mathbf{x}}
\newcommand{\dbf}{\mathbf{d}}
\newcommand{\Rbf}{\mathbf{R}}
\newcommand{\Tbf}{\mathbf{T}}

\newcommand{\Cbf}{\mathbf{C}}
\newcommand{\Ibf}{\mathbf{I}}
\newcommand{\Pbf}{\mathbf{P}}
\newcommand{\Qbf}{\mathbf{Q}}
\newcommand{\Vbf}{\mathbf{V}}
\newcommand{\Jbf}{\mathbf{J}}
\newcommand{\Xbf}{\mathbf{X}}
\newcommand{\Abf}{\mathbf{A}}
\newcommand{\Kbf}{\mathbf{K}}
\newcommand{\Gammabf}{\boldsymbol{\Gamma}}
\newcommand{\nubf}{\boldsymbol{\nu}}
\newcommand{\xibf}{\boldsymbol{\xi}}
\newcommand{\Xibf}{\boldsymbol{\Xi}}
\newcommand{\Omegabf}{\boldsymbol{\Omega}}


\newcommand{\ubf}{\mathbf{u}}

\newcommand{\lth}{\ell{\text{th}}}
\newcommand{\ith}{i{\text{th}}}
\newcommand{\jth}{j{\text{th}}}
\newcommand{\kth}{k{\text{th}}}
\newcommand{\ip}[2]{\left<#1,~#2\right>}

\newcommand{\OMIT}[1]{}

%%%%%%%%%%%%%%%%%%%%%%%%%%%%%%%%%%%%%%%%%%%%%%%%%%%%%%%%%
%                End Macros
%%%%%%%%%%%%%%%%%%%%%%%%%%%%%%%%%%%%%%%%%%%%%%%%%%%%%%%%%
\title{Homework 2}
%
%
% correct bad hyphenation here
\hyphenation{op-tical net-works semi-conduc-tor}
%\markboth{Submitted to IEEE Transactions on Control Systems
%Technology as a Brief Paper}
%         {Submitted to IEEE Transactions on Control Systems
%Technology as a Brief Paper}

\begin{document}
\maketitle


\begin{enumerate}

\item Suppose
$$U=\{(x,x,y)\in\mathbb{R}^3: x, y \in \re\}$$
Find a subspace $W$ of $\re^3$ such that $\re^3=U\oplus W$.


 \item For each subspace in (a)-(c), (1) find a basis, and (2) state the dimension.

 \begin{itemize}

  \item[(a)]

  $$
\left\{\left[\begin{array}{c}
x+ 2y \\
2x-3y \\
-x
\end{array}\right]: x, y \text { in } \mathbb{R}\right\}
$$

  \item[(b)]
$$
\left\{\left[\begin{array}{c}
x+ 3y-z \\
4x+5y+3z \\
3x+6z \\
-x+7y-9z
\end{array}\right]: x, y, z \text { in } \mathbb{R}\right\}
$$
  \item[(c)]
$$
\{(x, y, z, w): x-4 y+3w=0\}
$$
 \end{itemize}


 % \item $V$ is a nonzero finite-dimensional vector spaces, and the vectors listed belong to $V$. Mark each statement True or False.
 % \begin{itemize}
 %  \item[a.]
 %  If there exists a set $\left\{\mathbf{v}_{1}, \ldots, \mathbf{v}_{p}\right\}$ that spans $V$, then dim $V \leq p $.
 %  \item[b.]
 %  If there exists a linearly independent set $\left\{\mathbf{v}_{1}, \ldots, \mathbf{v}_{p}\right\}$ in $V$, then then dim $V \geq p$.
 %  \item[c.]
 %  If dim $V=p$, then there exists a spanning set of $p+1$ vectors in $V$.
 %  \item[d.]
 %  If there exists a linearly dependent set $\left\{\mathbf{v}_{1}, \ldots, \mathbf{v}_{p}\right\}$ in $V$, then dim $V \leq p$.
 %  \item[e.]
 %  If every set of $p$ elements in $V$ fails to span $V$, then dim $V>p$.
 %  \item[f.]
 %  If $p \geq 2$ and dim $V=p$, then every set of $p-1$ nonzero vectors is linearly independent.

 % \end{itemize}
 
\item Prove that If \(V_1\) and \(V_2\) are subspaces of a finite-dimensional vector space, then
\(\dim(V_1 + V_2)=\dim V_1+\dim V_2 - \dim(V_1\cap V_2)\). 

\item Consider two ordered bases $E=\{v_1,v_2,v_3\}$ and $F=\{w_1,w_2,w_3\}$ for a vector space $V$, and suppose
\begin{align*}
v_1=-w_1+w_2+w_3, \quad v_2=w_2+3w_3, \quad v_3= 4w_1-2w_2
\end{align*}
\begin{itemize}
\item[(a)] Find the transition matrix $S$ from $E$ to $F$.
\item[(b)] Compute the coordinate vector $[v]_F$ for $v=2v_1+1v_2-v_3$.
\end{itemize}


 \item Consider two ordered bases $E=\{v_1,v_2,v_3\}$ and $F=\{w_1,w_2,w_3\}$ for $\mathbb{R}^3$, where
\begin{align*}
v_1=\left[\begin{array}{c}4 \\ 6 \\ 7\end{array}\right], \quad v_2=\left[\begin{array}{c}0 \\ 1 \\ 1\end{array}\right], \quad v_3=\left[\begin{array}{c}0 \\ 1 \\ 2\end{array}\right]\\
w_1=\left[\begin{array}{c}1 \\ 1 \\ 1\end{array}\right], \quad w_2=\left[\begin{array}{c}1 \\ 2 \\ 2\end{array}\right], \quad w_3=\left[\begin{array}{c}2 \\ 3 \\ 4\end{array}\right]
\end{align*}
\begin{itemize}
\item[(a)] Find the transition matrix $S_1$ from $E$ to $F$.
\item[(b)] Find the transition matrix $S_2$ from $F$ to $E$.
\item[(c)] Verify that $S_1S_2=S_2S_1=I_3$.
\item[(d)] Compute the coordinate vector $[v]_E$, where $[v]=\left[\begin{array}{c} 4 \\ 6\\ 8\end{array}\right]$, and then use $S_1$ or $S_2$ (decide by yourself) to compute $[v]_F$.
\item[(e)] Check your work by computing $[v]_F$ directly.
\end{itemize}


\item Let $A\in\mathbb{R}^{4\times 5}$ be
\begin{align*}
A=\left[\begin{array}{cccccc} 1 & -2 &  1 & 1 & 2 \\
                             -1 &  3 &  2 & 0 & -2   \\
                              0 &  1 &  3 & 1 & 4 \\
                              1 &  2 & 13 & 5 & 5
   \end{array}\right].
\end{align*}
 \begin{itemize}
 \item[(a)] Find the four subspaces of the matrix ($C(A)$, $C(A^T)$, $N(A)$, and $N(A^T)$, determine their bases and dimensions).
   \item[(b)]
   Write down the fundamental theorem of linear algebra: Part 1 and Part 2. And verify them by the answers in (a).
 \end{itemize}



  \item Let $A\in\mathbb{R}^{4\times 5}$ and let $R$ be the reduced row echelon form of $A$. If the first and fourth columns of $A$ are
  \begin{align*}
  a_1=\left[\begin{array}{c}4 \\ 1 \\-2 \\ 0\end{array}\right], \quad a_4= \left[\begin{array}{c}-3 \\ -3 \\-1 \\ -5\end{array}\right]
  \end{align*}
and
\begin{align*}
R=\left[\begin{array}{ccccc}
1 & 0 & -1 & 0 & 1\\
0 & 1& -2 &0 & 0\\
0 & 0 & 0 & 1 & 2\\
0 & 0 & 0 & 0 & 0
\end{array}\right],
\end{align*}
\begin{itemize}
  \item[(a)] find a basis for $N(A)$.
  \item[(b)] given that $x_0$ is a solution to $Ax=b$, where
  \begin{align*}
  b=\left[\begin{array}{c}-2 \\ -3 \\ 0 \\ 1\end{array}\right], \quad x_0= \left[\begin{array}{c}0 \\ 3 \\ 1 \\ -1 \\ 0 \end{array}\right],
  \end{align*}
  determine the remaining column vectors of $A$.
  \end{itemize}
  
  
% \item If \(P\) is the plane of vectors in \(\mathbb{R}^4\) satisfying \(x_1 + x_2 + x_3 + x_4 = 0\), write a basis for \(P^{\perp}\). Construct a matrix that has \(P\) as its nullspace.  

\item Let $x$ and $y$ be linearly independent vectors in $\mathbb{R}^n$ and let $S=\operatorname{Span}(x,y)$. We can use $x$ and $y$ to define a matrix $A$ be setting
\begin{align*}
A=xy^T+yx^T
\end{align*}
\begin{itemize}
  \item[(a)] Show that $N(A)=S^{\perp}$.
  \item[(b)] Show that $\operatorname{dim} C(A)=2$ (the rank of $A$ must be $2$).
\end{itemize}
\end{enumerate}

%\bibliographystyle{IEEEtran}
%\bibliography{../../refs}
\end{document}